\section*{Capítulo \ref{ch:g11:absorbance} --- Cor e absorbância}

\begin{solution}{exo:g11:absorbance:1}
Absorvendo o laranja, ela parece azul, a cor diante do laranja no
círculo. Absorvendo o violeta, ela parece amarela.
\end{solution}

\begin{solution}{exo:g11:absorbance:2}
Uma solução verde absorve principalmente a luz vermelha, a cor diante do
verde no círculo: de cerca de \qty{620}{nm} a \qty{750}{nm}.
\end{solution}

\begin{solution}{exo:g11:absorbance:3}
Corante azul: $\lambda_{\max} = \qty{629}{nm}$, $A = 0.82$. Corante
amarelo: $\lambda_{\max} = \qty{427}{nm}$, $A = 0.795$, cerca de 0.80.
\end{solution}

\begin{solution}{exo:g11:absorbance:4}
$A$ é proporcional à concentração: o triplo dá $A = 0.90$; a metade dá
$A = 0.15$.
\end{solution}

\begin{solution}{exo:g11:absorbance:5}
O branco é uma cubeta cheia do \omterm{def:g6:solutions-and-solubility:solute}{solvente} puro. Ajustar $A = 0$ com ele
elimina o que a cubeta, o \omterm{def:g6:solutions-and-solubility:solute}{solvente} e o instrumento absorvem, de modo que
as leituras seguintes se devem só ao \omterm{def:g6:solutions-and-solubility:solute}{soluto}.
\end{solution}

\begin{solution}{exo:g11:absorbance:6}
$C = 0.41 / 0.164 = \qty{2.5}{mg/L}$.
\end{solution}

\begin{solution}{exo:g11:absorbance:7}
A \qty{629}{nm}, a \omterm{def:g11:absorbance:absorbance}{absorbância} é máxima, por isso a medida é mais
sensível (pequenas concentrações ainda dão \omterm{def:g11:absorbance:absorbance}{absorbâncias} legíveis); e no
alto da banda a curva é plana, por isso um pequeno erro no comprimento
de onda quase não muda a leitura. A \qty{550}{nm}, a \omterm{def:g11:absorbance:absorbance}{absorbância} é
pequena e varia depressa com o comprimento de onda.
\end{solution}

\begin{solution}{exo:g11:absorbance:8}
Diluída: $C = 0.48 / 0.164 = \qty{2.9}{mg/L}$; a bebida:
$5 \times 2.9 = \qty{15}{mg/L}$ (mais precisamente $14.6$).
\end{solution}

\begin{solution}{exo:g11:absorbance:9}
$a = A / (\ell\, C_m) = 0.795 / (1.00 \times 0.0150) = \qty{53.0}{L/(g.cm)}$;
$\varepsilon = a \times M = 53.0 \times 534.4 = \qty{2.83e4}{L/(mol.cm)}$.
\end{solution}

\begin{solution}{exo:g11:absorbance:10}
Acima de cerca de \qty{520}{nm}. A \qty{629}{nm}, o corante amarelo não
absorve nada, por isso toda a \omterm{def:g11:absorbance:absorbance}{absorbância} nesse ponto vem do corante
azul.
\end{solution}

\begin{solution}{exo:g11:absorbance:11}
Ela dobra: $A = \varepsilon \ell c$ é proporcional a $\ell$; a luz
atravessa o dobro de solução, e portanto o dobro de \omterm{def:g7:atoms-and-molecules:molecule}{moléculas}
absorventes.
\end{solution}

\begin{solution}{exo:g11:absorbance:12}
$A/C$: $0.164$, $0.162$, $0.115$, $0.0675$. A razão só é constante para
os dois primeiros: acima de cerca de \qty{10}{mg/L} (\omterm{def:g11:absorbance:absorbance}{absorbâncias} acima
de cerca de 1.6), a lei deixa de valer. Só os padrões de 5 e
\qty{10}{mg/L} (e os mais baixos) podem ser usados; uma solução
desconhecida que lê 2.5 deve ser diluída, por exemplo cinco vezes, e
medida de novo.
\end{solution}

\begin{solution}{exo:g11:absorbance:13}
Corante azul, só a \qty{629}{nm}: $C = 0.574 / 0.164 = \qty{3.5}{mg/L}$.
Corante amarelo, a \qty{427}{nm}: $C = 0.53 / 0.0530 = \qty{10}{mg/L}$.
\end{solution}

\begin{solution}{exo:g11:absorbance:14}
Ela encontra $0.368 / 0.164 = \qty{2.24}{mg/L}$. A \omterm{def:g11:absorbance:absorbance}{absorbância} verdadeira
é $0.368 - 0.040 = 0.328$, logo a concentração verdadeira é
$0.328 / 0.164 = \qty{2.00}{mg/L}$. Ela erra \qty{12}{\percent} para
mais.
\end{solution}

\begin{solution}{exo:g11:absorbance:15}
$10 \times 60 = \qty{600}{mg}$ por dia, em $600 / 20 = \qty{30}{L}$ de
bebida. Ninguém bebe trinta litros por dia: a bebida sozinha não
representa um risco realista, embora o corante também possa vir de
outros alimentos.
\end{solution}

\begin{solution}{pb:g11:absorbance:1}
\textbf{1.} A vermelho-alaranjada, a cor diante do azul no círculo.

\textbf{2.} $\lambda_{\max} = \qty{629}{nm}$.

\textbf{3.} A \qty{450}{nm}, a luz é azul: o corante a deixa passar, e é
exatamente por isso que a bebida parece azul.

\textbf{4.} A \qty{629}{nm}.

\textbf{5.} $\qty{50.0}{mg} / \qty{0.5000}{L} = \qty{100}{mg/L}$.

\textbf{6.} \omterm{def:g10:concentration-and-dilution:dilution}{Fator de diluição} $100 / 3.0 = 33.3$, logo
$V_0 = 100.0 / 33.3 = \qty{3.00}{mL}$: uma pipeta graduada (ou uma
bureta) e um balão volumétrico de \qty{100.0}{mL}.

\textbf{7.} $A/C$ = 0.168, 0.163, 0.165, 0.163, 0.165: todos perto de
0.164, logo $A \approx 0.164\, C$, uma reta que passa pela origem.

\textbf{8.} Uma solução sem corante ($C = 0$) é o branco, ajustado para
$A = 0$.

\textbf{9.} $a = \qty{0.164}{L/mg}$ por centímetro, isto é,
\qty{164}{L/(g.cm)}; $\varepsilon = 164 \times 792.9 =
\qty{1.30e5}{L/(mol.cm)}$.

\textbf{10.} $F = 50.0 / 25.0 = 2$.

\textbf{11.} $C = 0.590 / 0.164 = \qty{3.60}{mg/L}$.

\textbf{12.} $2 \times 3.60 = \qty{7.20}{mg/L}$.

\textbf{13.} Cerca de $2 \times 0.590 = 1.18$, acima do padrão mais alto
(0.823): a leitura ficaria fora da faixa calibrada, onde a lei pode
falhar. A \omterm{def:g10:concentration-and-dilution:dilution}{diluição} a traz para o meio dos padrões.

\textbf{14.} $\qty{7.20}{mg/L} \times \qty{0.500}{L} = \qty{3.60}{mg}$.

\textbf{15.} $n = \num{3.60e-3} / 792.9 = \qty{4.54e-6}{mol}$.

\textbf{16.} $6 \times 30 = \qty{180}{mg}$ por dia.

\textbf{17.} $180 / 3.60 = 50$ garrafas.

\textbf{18.} $50 \times 0.500 = \qty{25}{L}$ num dia: impossível de beber.
O corante desta bebida sozinho não pode levar uma criança perto do
limite, embora os corantes de vários alimentos se somem.

\textbf{19.} $6 \times 70 = \qty{420}{mg}$, isto é, $420 / 3.60 \approx
117$ garrafas.

\textbf{20.} Cerca de 50 garrafas de \qty{500}{mL}, \qty{25}{L} de
bebida, num único dia.
\end{solution}
